%%%%%%%%%%%%%%%%%%%%%%%%%%%%%%%%%%%%%%%%%%%%%%%%%%%%%%
% Plantilla para presentaciones
% en la Facultad de CC Matemáticas de la
% Universidad Complutense de Madrid.
%
% Autor: Ángel González Prieto
% Fecha: Noviembre 2021.
%%%%%%%%%%%%%%%%%%%%%%%%%%%%%%%%%%%%%%%%%%%%%%%%%%%%%%

\documentclass[11pt]{beamer}

\usetheme{Warsaw}
\usefonttheme{structurebold}
% Malmoe, Warsaw
\usepackage[utf8]{inputenc}
\usepackage[spanish,es-lcroman]{babel}
\usepackage{amsfonts,amsmath,amsthm}
\usepackage{enumitem}
% Remove navigation bar. Comment to show navigation bar.
\beamertemplatenavigationsymbolsempty

%%%%%%%%%%%%%%%%%%%%%% MATHEMATICAL COMMANDS %%%%%%%%%%%%%%%%%%%%%%%%%%%%

% Mathbold
\renewcommand{\AA}{\mathbb{A}}
\newcommand{\BB}{\mathbb{B}}
\newcommand{\CC}{\mathbb{C}}
\newcommand{\DD}{\mathbb{D}}
\newcommand{\FF}{\mathbb{F}}
\newcommand{\GG}{\mathbb{G}}
\newcommand{\HH}{\mathbb{H}}
\newcommand{\KK}{\mathbb{K}}
\newcommand{\NN}{\mathbb{N}}
\newcommand{\OO}{\mathbb{O}}
\newcommand{\PP}{\mathbb{P}}
\newcommand{\RR}{\mathbb{R}}
\newcommand{\TT}{\mathbb{T}}
\newcommand{\ZZ}{\mathbb{Z}}

% Calligraphic
\newcommand{\cA}{\mathcal{A}}
\newcommand{\cC}{\mathcal{C}}
\newcommand{\cD}{\mathcal{D}}
\newcommand{\cE}{\mathcal{E}}
\newcommand{\cF}{\mathcal{F}}
\newcommand{\cG}{\mathcal{G}} 
\newcommand{\cI}{\mathcal{I}} 
\newcommand{\cO}{\mathcal{O}} 
\newcommand{\cM}{\mathcal{M}} 
\newcommand{\cN}{\mathcal{N}} 
\newcommand{\cP}{\mathcal{P}} 
\newcommand{\cQ}{\mathcal{Q}} 
\newcommand{\cS}{\mathcal{S}} 
\newcommand{\cU}{\mathcal{U}} 
\newcommand{\cJ}{\mathcal{J}}
\newcommand{\cX}{\mathcal{X}}
\newcommand{\cY}{\mathcal{Y}}
\newcommand{\cZ}{\mathcal{Z}}
\newcommand{\cT}{\mathcal{T}}
\newcommand{\cV}{\mathcal{V}}
\newcommand{\cW}{\mathcal{W}}
\newcommand{\cB}{\mathcal{B}}
\newcommand{\cR}{\mathcal{R}}
\newcommand{\cH}{\mathcal{H}}

% Lie groups
\newcommand{\PGL}{\textup{PGL}}
\newcommand{\AGL}{\textup{AGL}}
\newcommand{\GL}{\textup{GL}}
\newcommand{\SL}{\textup{SL}}
\renewcommand{\O}{\textup{O}}
\newcommand{\U}{\textup{U}}
\newcommand{\SO}{\textup{SO}}
\newcommand{\SU}{\textup{SU}}
\newcommand{\norm}[1]{\lVert#1\rVert} %norma%


%%%%%%%%%%%%%%%%%%%%%%%%%%%
%                   FRONT SLIDE INFO                            %
%%%%%%%%%%%%%%%%%%%%%%%%%%%


\title{Introducción a la Teoría de Morse}

\author{Enrique Díaz Blanco}

\institute[Universidad Complutense de Madrid]
{Universidad Complutense de Madrid\\[0.3cm]
	\begin{figure}[h]
	\begin{center}
	
\includegraphics[scale=0.15]{/home/kike/Documents/Mates/TFG/introduction-to-morse-theory/presentacion/imagenes/logocomplu.png}
	\end{center}
	\end{figure}
}

\date{}

%%%%%%%%%%%%%%%%%%%%%%%%%%%%%%%%%%%%%%%%%%%%%%%%%%%%%%%%%%%%%%%%%%%%%%%%%%%%%%%%%

\begin{document}

% Theorem-like environments
\newtheorem*{thm}{Teorema}
\newtheorem*{prop}{Proposición}
\newtheorem*{lem}{Lema}
\newtheorem*{cor}{Corolario}

\theoremstyle{definition}
\newtheorem*{defn}{Definición}
\newtheorem*{ex}{Ejemplo}
\newtheorem*{rmk}{Observación}
\newtheorem*{conj}{Conjetura}


%%%%%%%%%%%%%%%%%%%%%%%%%%%
%                            SLIDES                                       %
%%%%%%%%%%%%%%%%%%%%%%%%%%%

\frame{\titlepage} % Show title slide
\part{Main Part}

%%%%%% FLOWCHART %%%%%%%




%%%%%%%%%%%%%%%%%%%%%%%%%%%%%%%%%%%%
%%                                %%
%%           Capitulo 1           %%  
%%                                %%
%%%%%%%%%%%%%%%%%%%%%%%%%%%%%%%%%%%%

\section{Puntos críticos}
\begin{frame}
\frametitle{Definición de punto crítico}
	\begin{definition}
    Sea $f: $M$ \to \mathbb{R}$ una función suave.
    \begin{enumerate}[label=(\roman*)]
        \item Diremos que $p\in \mathrm{M}$ es un \textit{punto crítico de $f$} si $d_{p}f:T_{p}M\to \mathbb{R}$ es idénticamente cero. Es decir, dado un sistema de coordenadas local $(x_{1}, \dots ,x_{n})$ en un entorno $\mathcal{U}$ de $p$ tenemos
        \begin{equation*}
            \frac{\partial f}{\partial x_1}(p)= \dots = \frac{\partial f}{\partial x_n}(p) = 0.
        \end{equation*}
        
        \item A su vez, diremos que $f(p)\in \mathbb{R}$ es un \textit{valor crítico de $f$} si $p \in \mathrm{M}$ es un punto crítico.
        
    \end{enumerate}
	\end{definition}
\end{frame}



% Slide

\begin{frame}
\frametitle{Índice de un punto crítico}
	\begin{definition}
    	Sea $p\in M$ un punto crítico de f. Entonces $p$ es \textit{no degenerado} si la matriz  $\left(\frac{\partial^{2}f}{\partial x_i \partial x_j}(p)\right)$ es no singular.
	\end{definition}
	\bigskip
	En coordenadas, $\text{Hess}_pf = \left(\frac{\partial^{2}f}{\partial x_i \partial x_j}(p)\right)$ 
	\bigskip
	\begin{definition}
		Llamamos índice de $p$ a la máxima dimensión de un subespacio de $T_pM$ sobre el que $\text{Hess}_pf$ es definida negativa.
	\end{definition}
\end{frame}

\begin{frame}
\frametitle{El Lema de Morse}

\begin{alertblock}{Lema de Morse}
    Sean $p$ un punto crítico \textbf{no degenerado} de una función suave $f:M\to \mathbb{R}$. Entonces existe un sistema de coordenadas local $(y_1,\dots, y_n)$ en un entorno $\mathcal{U}$ de $p$ tal que $y_i(p)=0, \forall i\in \{1\dots,n\}$ y se verifica
    \begin{equation*}
        f=f(p)-(y_1)^2-\dots-(y_\lambda)^2 + (y_{\lambda +1})^2+\dots + (y_n)^2
    \end{equation*}
    en todo $\mathcal{U}$ con $\lambda$ el índice de $f$ en $p$.
\end{alertblock}

\begin{corollary}
	Todo punto crítico no degenerado de una función suave de $M$ en $\mathbb{R}$ está \textbf{aislado} respecto de los demás puntos críticos.
\end{corollary}

\end{frame}


\begin{frame}
	\begin{definition}
		Diremos que una función suave $f: M\to \mathbb{R}$ es una \textit{función de Morse} si todos sus puntos críticos son no degenerados.
	\end{definition}

\end{frame}

\begin{frame}
	\begin{example}
		$f(x,y)=(x^2-1)^2 + (y^2-1)^2$
	\end{example}
	
	\begin{figure}[h]
	\begin{center}
	\includegraphics[scale=0.45]{/home/kike/Documents/Mates/TFG/introduction-to-morse-theory/presentacion/imagenes/grafica1.png}
	\end{center}
	\end{figure}
\end{frame}




%%%%%%%%%%%%%%%%%%%%%%%%%%%%%%%%%%%%
%%                                %%
%%           Capitulo 2           %%  
%%                                %%
%%%%%%%%%%%%%%%%%%%%%%%%%%%%%%%%%%%%

\section{CW-complejos y tipo de homotopía}

\begin{frame}
\frametitle{Tipo de homotopía}

Denotamos la relación de homotopía por el símbolo $\simeq$.
\bigskip
\begin{definition}
	Dos espacios $X$ e $Y$ tienen el mismo \textbf{tipo de homotopía} si existe una equivalencia homotópica entre ellos.
\end{definition}

\bigskip

\begin{definition}
	$f:X\to Y$ es una equivalencia homotópica si existe $g:Y\to X$ tal que
	\begin{equation*}
    g \circ f \simeq \mathrm{id}_X \qquad \text{y} \qquad f \circ g \simeq \mathrm{id}_Y.
	\end{equation*}

\end{definition}

\end{frame}

\begin{frame}

	\begin{definition}
		Sea $X$ un espacio topológico y $A \subset X$ un subespacio. Decimos que $A$ es un
		\textbf{retracto de deformación fuerte} de $X$ si existe una homotopía relativa a $A$ de la función identidad de $X$, $id_X$, a una retracción $r: X \to A$.
	\end{definition}

	\begin{alertblock}{}
		\textbf{$A$ tiene el mismo tipo de homotopía que $X$.}
	\end{alertblock}

\end{frame}

\begin{frame}
\frametitle{CW-complejos}
\begin{definition}
    Llamamos \textbf{CW-complejo} a un espacio topológico construido siguiendo los siguientes pasos:
	\vspace{1em}
        \begin{enumerate}[label=(\roman*), itemsep=2em]
            \item Se comienza con un conjunto discreto $ X_0 $, el \textbf{$0$-esqueleto}.
            
            \item Se construye, de \textbf{manera inductiva}, el $ n $-esqueleto para $n>0$ como
            \begin{equation*}
                X_n = X_{n-1}\cup_{\varphi_{\alpha}}D^n_{\alpha},~\alpha\in \Lambda_n.
            \end{equation*}
            
            \item Llamaremos \textit{CW-complejo} al espacio 
            \begin{equation*}
                X = \bigcup_{n} X_n.    
            \end{equation*}
        \end{enumerate} 
\end{definition}
\end{frame}

\begin{frame}
\frametitle{El toro plano como CW-complejo}
	\begin{columns}
	
	\column{0.4\textwidth}

	Se puede construir como un CW-complejo con:
	\vspace{1em}
	\begin{itemize}
		\item Una $0$-celda.
		\item Dos $1$-celdas.
		\item Una $2$-celda.
	\end{itemize}

	\column{0.6\textwidth}

	\begin{figure}[h]
	\begin{center}
	\includegraphics[scale=0.4]{/home/kike/Documents/Mates/TFG/introduction-to-morse-theory/presentacion/imagenes/toroplano.pdf}
	\end{center}
	\end{figure}

	\begin{figure}[h]
	\begin{center}
	\includegraphics[scale=0.4]{/home/kike/Documents/Mates/TFG/introduction-to-morse-theory/presentacion/imagenes/frontera_toroplano.pdf}
	\end{center}
	\end{figure}
	\end{columns}
	
\end{frame}


%%%%%%%%%%%%%%%%%%%%%%%%%%%%%%%%%%%%

\section{Sobre el pegado de $\lambda$-celdas}

\begin{frame}
	
\end{frame}


\end{document}
